% Generated draft: requires mathematical and source review.
\documentclass[11pt]{article}
\usepackage[margin=25mm]{geometry}
\usepackage{fontspec}
\setmainfont{Latin Modern Roman}
\usepackage{amsmath,amssymb,mathtools,unicode-math}
\setmathfont{Latin Modern Math}
\usepackage{enumitem,needspace,xurl,hyperref}
\setlength{\parindent}{0pt}
\setlength{\parskip}{5pt}
\setlength{\emergencystretch}{3em}
\newcommand{\N}{\mathbb{N}}
\newcommand{\Z}{\mathbb{Z}}
\newcommand{\Q}{\mathbb{Q}}
\newcommand{\R}{\mathbb{R}}
\newcommand{\C}{\mathbb{C}}
\newcommand{\F}{\mathbb{F}}
\newcommand{\A}{\mathbb{A}}
\newcommand{\PP}{\mathbb{P}}
\newcommand{\E}{\mathbb{E}}
\newcommand{\eps}{\varepsilon}
\newcommand{\dd}{\,\mathrm{d}}
\begin{document}
\section*{UnsolvedMath ID: 701091}
\textbf{Source problem code:} catalogue entry
\subsection{Definitions and mathematical statement}
For a finite group \(G\), let \(\operatorname{Irr}(G)\) denote its ordinary irreducible complex characters; this uses the customary interpretation of the unqualified word ``characters'' in the supplied record. For \(\chi\in\operatorname{Irr}(G)\), its degree is \(\chi(1)\), and its kernel is \(\ker\chi=\{g\in G:\chi(g)=\chi(1)\}\). The character is faithful if \(\ker\chi=\{1\}\), and non-linear if \(\chi(1)>1\). Define the pointwise product by \((\alpha\beta)(g)=\alpha(g)\beta(g)\) for every \(g\in G\). A group is solvable if its derived series, given by \(G^{(0)}=G\) and \(G^{(i+1)}=[G^{(i)},G^{(i)}]\), reaches the trivial group after finitely many steps.

\textbf{Problem (I. M. Isaacs, Kourovka Notebook 15.3).} Does there exist a finite solvable group \(G\) and characters \(\alpha,\beta\in\operatorname{Irr}(G)\) such that
\[
\ker\alpha=\ker\beta=\{1\},\qquad \alpha(1)>1,\qquad \beta(1)>1,\qquad \alpha\beta\in\operatorname{Irr}(G)?
\]
No requirement that \(\alpha\ne\beta\) is stated in the source.

The supplied record states that examples exist for non-solvable groups, including \(\mathrm{SL}_2(5)\) with two irreducible characters of degree \(2\). It further attributes to I. Zisser the result that the alternating group \(\mathbb{A}_n\) admits such characters if and only if \(n\) is a square exceeding \(4\). The alternating-group notation repairs an apparent typesetting error in the supplied text. The record describes the cited work of I. M. Isaacs as giving evidence, but no proof, that solvable examples do not exist. These contextual claims and the supplied \texttt{open} label are source reports, not independently verified assessments.
\subsection{Short English statement}
Can a finite solvable group have two faithful irreducible complex characters, each of degree greater than \(1\), whose pointwise product is irreducible? The supplied source reports non-solvable examples and evidence, but no proof, against solvable examples.
\subsection{Sources}
\begin{itemize}
\item Local UnsolvedMath registry, supplied complete source record: permanent ID \(701091\), exact source problem code \texttt{catalogue entry}, title ``Kourovka Notebook Problem 15.3.'' The record attributes the problem to I. M. Isaacs, proposed in 2002, and labels it \texttt{open}; this label has not been independently verified.
\item \emph{The Kourovka Notebook (Unsolved Problems in Group Theory)}, original problem 15.3, 15th issue (2002); the supplied background also identifies the 21st issue edition (2026), arXiv:1401.0300v48. The supplied collection description names E. I. Khukhro and V. D. Mazurov as editors.
\item I. Zisser, \emph{Israel J. Math.}, \textbf{84}, nos. 1--2 (1993), 147--151. Cited in the supplied statement for the alternating-group result; no article title or URL is supplied.
\item I. M. Isaacs, \emph{J. Algebra}, \textbf{223}, no. 2 (2000), 630--646. Cited in the supplied statement as providing evidence, but no proof, against solvable examples; no article title or URL is supplied.
\end{itemize}
\end{document}
